% updated Jan 2026 by Dim Papadopoulos
% updated September 2023 by Alessio Del Bue
% updated April 2002 by Antje Endemann
% Based on CVPR 07 and LNCS, with modifications by DAF, AZ and elle, 2008 and AA, 2010, and CC, 2011; TT, 2014; AAS, 2016; AAS, 2020; TH, 2022

\documentclass[runningheads]{llncs}
\usepackage{graphicx}
% DO NOT USE \usepackage{times}, it will be removed by typesetters
%\usepackage{times}

\usepackage{tikz}
\usepackage{comment,verbatim}
\usepackage{amsmath,amssymb} % define this before the line numbering.
\usepackage{color}

% The "axessiblity" package can be found at: https://ctan.org/pkg/axessibility?lang=en
\usepackage[accsupp]{axessibility}  % Improves PDF readability for those with disabilities.


% INITIAL SUBMISSION - The following two lines are NOT commented
% CAMERA READY - Comment OUT the following two lines
\usepackage{ruler}
\usepackage[width=122mm,left=12mm,paperwidth=146mm,height=193mm,top=12mm,paperheight=217mm]{geometry}



\begin{document}
% \renewcommand\thelinenumber{\color[rgb]{0.2,0.5,0.8}\normalfont\sffamily\scriptsize\arabic{linenumber}\color[rgb]{0,0,0}}
% \renewcommand\makeLineNumber {\hss\thelinenumber\ \hspace{6mm} \rlap{\hskip\textwidth\ \hspace{6.5mm}\thelinenumber}}
% \linenumbers
\pagestyle{headings}
\mainmatter
%\def\ECCVSubNumber{100}  % Insert your submission number here

\title{Workshop Proposal Template for ECCV 2026} % Replace with your title

% INITIAL SUBMISSION 
%\begin{comment}
\titlerunning{Workshop Proposal Template for ECCV 2026} 

\author{First Organizer\inst{1} \and
Second Organizer\inst{2,3} \and
Third Organizer\inst{3}}
%
\authorrunning{Workshop Proposal Template for ECCV 2026} 
% First names are abbreviated in the running head.
% If there are more than two authors, 'et al.' is used.
%
\institute{Princeton University, Princeton NJ 08544, USA \and
Springer Heidelberg, Tiergartenstr. 17, 69121 Heidelberg, Germany
\email{lncs@springer.com}\\
\url{http://www.springer.com/gp/computer-science/lncs} \and
ABC Institute, Rupert-Karls-University Heidelberg, Heidelberg, Germany\\
\email{\{abc,lncs\}@uni-heidelberg.de}}
%\end{comment}
%******************



%\author{Anonymous ECCV submission}
%\institute{Paper ID \ECCVSubNumber}
%\end{comment}
%******************

% CAMERA READY SUBMISSION
\begin{comment}
\titlerunning{Abbreviated paper title}
% If the paper title is too long for the running head, you can set
% an abbreviated paper title here
%
\author{First Author\inst{1}\orcidID{0000-1111-2222-3333} \and
Second Author\inst{2,3}\orcidID{1111-2222-3333-4444} \and
Third Author\inst{3}\orcidID{2222--3333-4444-5555}}
%
\authorrunning{F. Author et al.}
% First names are abbreviated in the running head.
% If there are more than two authors, 'et al.' is used.
%
\institute{Princeton University, Princeton NJ 08544, USA \and
Springer Heidelberg, Tiergartenstr. 17, 69121 Heidelberg, Germany
\email{lncs@springer.com}\\
\url{http://www.springer.com/gp/computer-science/lncs} \and
ABC Institute, Rupert-Karls-University Heidelberg, Heidelberg, Germany\\
\email{\{abc,lncs\}@uni-heidelberg.de}}
\end{comment}
%******************
\maketitle

%\begin{abstract}
%The abstract should summarize the contents of the paper. LNCS guidelines
%indicate it should be at least 70 and at most 150 words. It should be set in 9-point
%font size and should be inset 1.0~cm from the right and left margins.
%\dots
\keywords{We would like to encourage you to list here your keywords related to the Workshop/Tutorial topic.}
%\end{abstract}


\section{Workshop Submission Guidelines}


Workshop proposals should be in PDF format and limited to 5 pages using this template (references not counted for page limit). They should be submitted to the Chairs at \url{https://forms.gle/9gHLi2kHtSFZsF3V6} by 27th of February 2026, including the following information:
\begin{itemize}
    \item Workshop title (use \verb!\title{}!).
    \item Proposers’ names, titles, affiliations, and primary contact email (use \verb!\author{}!, \verb!\institute{}!, \verb!\email{}!).
    \item Relevant keywords summarizing tutorial topic (use \verb!\keywords{}!).
    \item Relevance of Topic: Description of the topics that will be covered and their relevance to the computer vision community.
    \item Relevance of Organizers: Experience that makes the proposers well-suited for organizing the workshop and their relevance in the computer vision community.
    \item Tentative program outline (including preference for a half- or full-day event, estimated numbers of orals, posters, and invited talks)
    \item Names and bios of any tentative/confirmed invited speakers
    \item How diversity and inclusion in the organizing team and speakers is promoted (gender, race, seniority, affiliations, interdisciplinarity of topics and points of view).
    \item Anticipated target audience as well as expected number of attendees.
    \item Description of how this proposal relates to previous workshops at CVPR, ICCV, ECCV (be as specific as possible).
    \item Any special space or equipment requests.
\end{itemize}

Note: two special sections must be included: Relevance of Topic (max 2000 characters) and Relevance of Organizers (max 1500 characters). These must have the same content as the corresponding fields of the submission form.

For any questions, please contact the workshop chairs via email: \\ \email{eccv2026-workshopchairs@ecva.net}.


% \clearpage

% ---- Bibliography ----
%
% BibTeX users should specify bibliography style 'splncs04'.
% References will then be sorted and formatted in the correct style.
%
% \bibliographystyle{splncs04}
% \bibliography{egbib}
\end{document}
